
\chapter{Metodología de ingeniería de software}
\label{chap:engineering-methodology}

\section{Enfoque de desarrollo}

El desarrollo de APEX sigue un enfoque incremental e iterativo,
en el que las funcionalidades se incorporan progresivamente
a partir de componentes independientes.

Este enfoque permite dividir el sistema en unidades de
responsabilidad bien definidas, facilitando su implementación,
verificación y evolución.

En lugar de desarrollar simultáneamente todas las capacidades
previstas para la plataforma, se prioriza la construcción
de una infraestructura computacional sobre la cual puedan
incorporarse posteriormente los componentes científicos.

Cada componente desarrollado debe contar con una
responsabilidad identificable, interfaces definidas y
mecanismos que permitan verificar su comportamiento.

El proceso de desarrollo contempla las siguientes actividades:

\begin{enumerate}
    \item Identificación de los requisitos del componente.
    \item Diseño de su estructura y responsabilidades.
    \item Definición de los casos de prueba.
    \item Implementación de la funcionalidad.
    \item Ejecución de pruebas y análisis estático.
    \item Corrección de defectos y refinamiento del diseño.
    \item Integración con el resto de la infraestructura.
\end{enumerate}

\section{Arquitectura modular}

La arquitectura de APEX se organiza en módulos que agrupan
funcionalidades relacionadas.

Esta organización busca reducir el acoplamiento entre
componentes y favorecer la cohesión interna de cada módulo.

La infraestructura computacional inicial se divide
conceptualmente en las siguientes áreas:

\begin{itemize}
    \item \textbf{Hardware:} detección de las características
    y recursos del sistema.

    \item \textbf{Execution:} ejecución de tareas y
    administración de trabajadores.

    \item \textbf{Logging:} registro estructurado de eventos
    y operaciones del sistema.
\end{itemize}

La separación de responsabilidades permite que los
componentes puedan evolucionar de manera independiente
siempre que se mantengan sus contratos e interfaces.

Asimismo, proporciona una base para incorporar en el futuro
subsistemas de simulación, modelado físico y optimización
computacional.

\section{Desarrollo guiado por pruebas}

Durante el desarrollo de APEX se utiliza el enfoque
Test-Driven Development (TDD), mediante el cual las pruebas
se diseñan antes de implementar la funcionalidad
correspondiente.

El ciclo de desarrollo se basa en tres etapas:

\begin{enumerate}
    \item \textbf{Red:} escribir una prueba que describa
    un comportamiento esperado y comprobar que inicialmente
    falla por la ausencia de la implementación.

    \item \textbf{Green:} desarrollar la implementación
    mínima necesaria para satisfacer la prueba.

    \item \textbf{Refactor:} mejorar la estructura interna
    del código sin alterar el comportamiento verificado.
\end{enumerate}

Este procedimiento permite establecer expectativas
explícitas sobre el comportamiento de los componentes
y detectar regresiones durante su evolución.

La suite de pruebas se desarrolla utilizando
\texttt{pytest}, con pruebas unitarias organizadas
por componente.

\section{Calidad y análisis estático}

La calidad del código se evalúa mediante herramientas
complementarias que analizan diferentes propiedades
del software.

En el estado actual del proyecto se utilizan:

\begin{itemize}
    \item \textbf{pytest:} ejecución de pruebas automatizadas.
    \item \textbf{Ruff:} análisis estático y aplicación
    consistente de convenciones de estilo.
    \item \textbf{Mypy:} comprobación estática de tipos.
    \item \textbf{Ruff Formatter:} normalización del formato
    del código fuente.
\end{itemize}

Estas herramientas cubren aspectos diferentes
del proceso de verificación. La ejecución satisfactoria
de las pruebas no sustituye al análisis estático,
y ninguno de estos mecanismos demuestra por sí solo
la corrección científica de los modelos físicos.

Por esta razón, la validación de los componentes
computacionales y la validación de los modelos matemáticos
se consideran actividades distintas.

\section{Gestión del código fuente}

El desarrollo utiliza Git como sistema de control de versiones
y GitHub como plataforma de alojamiento del repositorio.

Se emplean ramas de desarrollo para organizar la
incorporación de funcionalidades y mantener un historial
de cambios identificable.

Los commits permiten asociar modificaciones concretas
con las funcionalidades implementadas.

Esta práctica contribuye a la trazabilidad del desarrollo
y facilita la identificación de cambios que puedan afectar
el comportamiento del sistema.

\section{Reproducibilidad y validación}

La reproducibilidad constituye un principio importante
para el desarrollo de una plataforma orientada a la
computación científica.

Los resultados computacionales deben poder relacionarse
con las condiciones bajo las cuales fueron obtenidos,
incluyendo las configuraciones utilizadas, los parámetros
de ejecución y las versiones relevantes del software.

En la etapa actual, la verificación de APEX se centra
principalmente en pruebas automatizadas de sus componentes
de infraestructura.

En etapas posteriores será necesario incorporar
procedimientos adicionales de validación numérica
y científica, incluyendo:

\begin{itemize}
    \item Comparación con soluciones analíticas cuando
    estén disponibles.
    \item Comparación con resultados de referencia
    publicados.
    \item Evaluación de convergencia numérica.
    \item Análisis de sensibilidad respecto a los parámetros
    del modelo.
    \item Repetición de experimentos bajo condiciones
    controladas.
\end{itemize}

Estos procedimientos serán especialmente relevantes
cuando se incorporen modelos físicos, simulaciones
y algoritmos de optimización.

\section{Consideraciones metodológicas}

El enfoque adoptado busca mantener una separación clara
entre el desarrollo del software, la verificación de su
implementación y la validación científica de los resultados.

Esta distinción evita interpretar la ausencia de errores
de programación como evidencia suficiente de que un modelo
matemático representa correctamente un fenómeno físico.

Por tanto, APEX considera la calidad del software,
la reproducibilidad computacional y el rigor científico
como dimensiones complementarias del proceso de desarrollo.
